
\section{Closed forms and methods}
\label{sec:at-closed}

\begin{proposition}[Scaled Schwinger form of a defect cycle]
\label{prop:at-schwinger}
For a cycle of Lemma \ref{lem:at-special} with $n_t$ edges, external momenta $p_1,\dots,p_n$ and superficial degree $D$, substitute $t_e=\varepsilon\tau_e$, $\tau_\sigma\in[0,1]$ for the special edge and $\tau_e\in[1,\infty)$ for the others (both segments of a composite line included), and integrate the loop momentum against the Gaussian. With $\tau=\sum_e\tau_e$, $\alpha_e=\tau_e/\tau$, the coefficient of $\varepsilon^{j-D/2}$ of the Taylor expansion in the external momenta is the order-$j$ Taylor coefficient of the Gaussian-integrated integrand integrated over the scaled region; for $j<D$ odd orders vanish by Lorentz parity and even orders give the poles $a_{(D-j)/2}$, the order $j=D$ gives $a_0$, and the orders $j>D$ give $O(\varepsilon^{1/2}\log\varepsilon)$ or smaller. Each $a_0$ is a finite sum of master integrals
\[
J(n_t;\beta;N)=\int_0^1d\tau_\sigma\int_{[1,\infty)^{n_t-1}}\prod_{e\ne\sigma}d\tau_e\;\tau^\beta\,\tau^{-N},
\]
($\tau^\beta$ a monomial in the $\tau_e$), convergent when $N>|\beta|+n_t-1$, with values in $\mathbb Q+\mathbb Q\log2+\mathbb Q\log3$ for $n_t\le3$; on the simplex this is the formula of Lemma \ref{lem:at-special}.
\end{proposition}

\begin{proof}
Theorem c9-sing (i) and its proof, which applies to every cycle with one special edge; the master integrals are those of the Part-C engine of Step 2 (route A), evaluated exactly and confirmed at $50$ digits.
\end{proof}

\begin{proposition}[Small closed forms]
\label{prop:at-small}
With $F_m$ the $m$-th simplex integrand:
\begin{itemize}
\item (i) $n_t=2$: $16\pi^2a_0=\int_0^{1/2}F_0(\alpha)\log\frac{1-\alpha}\alpha\,d\alpha$, $\alpha=\alpha_\sigma$; for constant $F_0$ this is $F_0\log2$;
\item (ii) constant $F_0$ ($D=0$), any $n_t$: $16\pi^2a_0=F_0\,\frac{1}{(n_t-1)!}\log\frac{n_t}{n_t-1}$; constant $F_1$ ($D=2$): $16\pi^2a_1=\frac{F_1}{n_t^2(n_t-1)\,(n_t-1)!}$;
\item (iii) one vertex of $I$ and a composite line ($k=1$, Proposition c11-k1): $16\pi^2a_0=\int_0^{1/3}F_0(\beta)\big[(1-\beta)\log\frac{1-\beta}{2\beta}-(1-3\beta)\big]d\beta$, in $\mathbb Q+\mathbb Q\log2+\mathbb Q\log3$;
\item (iv) triangles of cubic vertices with the special edge: values in $\mathbb Q+\mathbb Q\log\tfrac32$ (the table of Section 9.1).
\end{itemize}
\end{proposition}

\begin{proof}
(i) For $n_t=2$ the region where $\alpha_\sigma$ is the smaller coordinate is $\alpha\le\tfrac12$ and $\min_{f\ne\sigma}\alpha_f=1-\alpha$; $\int_0^{1/2}\log\frac{1-\alpha}\alpha d\alpha=\log2$. (ii) By symmetry the region where $\alpha_\sigma$ is smallest is $1/n_t$ of the simplex, of volume $1/(n_t-1)!$; the logarithmic average over it is $\log\frac{n_t}{n_t-1}$ (Part Q, closed forms B2). (iii) Proposition c11-k1. (iv) Theorem c9-sing and the table of Section 9.1.
\end{proof}

\begin{proposition}[Position factors and Killing data]
\label{prop:at-inject}
For $x\in\{\Delta,\mathfrak{so}(V),\eta,K_b\}$, a vertex $S(x)$ at the point $y$ carries a polynomial of degree $p_x$ in $y$. With a momentum $r$ injected at that vertex and absorbed by the plane wave of the functional, $y^\mu\mapsto-i\partial/\partial r_\mu$ and $y^\mu y^\nu\mapsto-\partial^2/\partial r_\mu\partial r_\nu$ at $r=0$, acting on the integrand and on the momentum delta function. The action on the delta function gives the position parts of Proposition \ref{prop:at-pos}; the action on the integrand gives the intrinsic part for the chosen canonical densities (Definition c11-canonical: all momenta but one held fixed). The $\varepsilon^0$ coefficient is polynomial in $r$, so exact Lagrange interpolation on integer $r$ reproduces the derivatives.
\end{proposition}

\begin{proof}
$\int\mathrm{dvol}\,y^\mu e^{i(r,y)}(\cdots)=-i\partial_{r_\mu}\int\mathrm{dvol}\,e^{i(r,y)}(\cdots)$; Lemma \ref{lem:at-special} gives polynomiality.
\end{proof}

\begin{proposition}[Second route: renormalised amplitudes]
\label{prop:at-routeB}
For each generator and each sector, the renormalised arity-one amplitude $W^{(1)}_{\rm ren}(x)$ of Theorem \ref{thm:at-second} is a sum of one-loop integrals with propagators $e^{-\varepsilon|q|^2}|q|^{-2}$ (and $e^{-2\varepsilon|q|^2}|q|^{-4}$ times the numerator of $(Q^{GF})^2$ on composite lines), minus the trees of $\ICT$. Their finite parts at $\varepsilon\to0$ are Feynman-parameter integrals of rational functions times logarithms of the external invariants; $Q$ acts on the plane-wave value by replacing one leg by its image under the substitutions of Proposition c1-Q, and the nonlocal terms ($\log$ of momenta, ratios $(p,a)(p,b)|p|^{-2}$) cancel in $Q\,W^{(1)}_{\rm ren}(x)$ up to the trees of the components with fewer legs. A third evaluation of each special-edge cycle at finite $\varepsilon$ (numerical, $\varepsilon=2^{-k}$, Richardson extrapolation of $\varepsilon\,\mathcal C$ and $\mathcal C-\sum_m\varepsilon^{-m}a_m$) checks the absence of $\log\varepsilon$ and the value of $a_0$.
\end{proposition}

\begin{proof}
Theorem \ref{thm:at-second} (i); the Feynman-parameter form follows from the Schwinger form by the substitution $t_e=\tau\alpha_e$ and integration over $\tau$; Lemma \ref{lem:at-special} gives the expansion used in the extrapolation.
\end{proof}

\begin{remark}[Theorems of the construction used as cross-checks]
\label{rem:at-cross}
The following identities are consequences of the construction and are used only as cross-checks of the two routes: Theorem \ref{thm:at-second} (v), $\Ob^{(1)}(\Delta)=4\lc^{(0)}+(\text{composite part})$, whose two-leg form is $u[-4,16,-4]+u[\tfrac13,-\tfrac13,-\tfrac23]$ on $\int(dA,dA)$, $\int(\lambda,\dslash\lambda)$, $\int(dc,A^*)$ (X-scal); the value $\theta=-\tfrac13$ of Theorem c9-kappa (X-theta); the twisted obstruction of Remark \ref{rem:at-twist} against the twist literature (X-twist); and $K=P^\dagger$ (Corollary c8-gram) for the special conformal data (X-KP). None of them is used to produce a number.
\end{remark}

\section{Computed two-leg values, solutions and relations}
\label{sec:at-computed}

Values of the two engine routes of Section \ref{sec:at-engines} (step2/strict/action/engine\_a, engine\_b, complete), equal entry by entry in $\mathbb Q(i)+\mathbb Q(i)\log2+\mathbb Q(i)\log3$, with the signs of Lemma \ref{lem:at-sformulas}. Units $u$; $L:=\log\tfrac32$; $T_0,\dots,T_4$, $g_1,\dots,g_4$, $s_0$, $S(\Delta)_d$, $S(P_b)_d$, $N$ as in Proposition \ref{prop:at-xi2} and Remark \ref{rem:at-twocomp}; $u_0:=\int(\rho(\AV)\eta,\lambda^*)$; $S(M_\omega)_{\rm spin}$ as in Proposition \ref{prop:at-two} (iii); $F^{\rm part}(\xi'):=u(L-\tfrac{79}{12})(T_1+T_2)(\xi')$ for $\xi'\in\Sigma$.

\begin{proposition}[Arity-one obstructions on two legs]
\label{prop:at-c-ob1}
\[
\begin{aligned}
&\hbar\Ob^{(0)}|_2=-\tfrac{27}8u\int(d^*dc,d^*\AV),\qquad \hbar\Ob^{(1)}(\Delta)|_2=\eqref{eq:at-reg2},\qquad \hbar\Ob^{(1)}(P_b)|_2=0,\qquad \hbar\Ob^{(1)}(R_r)|_2=0,\\
&\hbar\Ob^{(1)}(M_\omega)|_2=\tfrac13u\,Q\,S(M_\omega)_{\rm spin},\qquad \hbar\Ob^{(1)}(\xi)|_2=u\Big[\big(\tfrac{79}{24}-\tfrac12L\big)(g_1+g_3)+\tfrac{27}8g_2+0\cdot g_4\Big],\\
&\hbar\Ob^{(1)}(\eta)|_2=\sum_\mu x^\mu\,\hbar\Ob^{(1)}(\rho(e_\mu)\eta)|_2+u\Big[-\tfrac52\int(\eta,\lambda)\,d^*\AV-\tfrac52\int(\eta,\rho(d\AV)\lambda)+\big(-\tfrac32+3L\big)\int(\eta,\rho(d\Ap)\lambda)+\tfrac52\int(\rho(dc)\eta,\lambda^*)\Big]+Q(\cdot),
\end{aligned}
\]
the intrinsic part of the last line in the canonical density (Definition c11-canonical); along $K_b$ the obstruction is in the span of the canonical data with position degree $\le2$, with pole sum $\tfrac94\{\int(A,A),S(K_b)\}$. $\lc^{(1)}(\xi)|_2=\{uN,S_2(\xi)\}=2u(T_1+T_2-T_3-T_4)$, so $l_\xi=0$ in Remark \ref{rem:at-logexact} on two legs. The Wess--Zumino identities of Lemma \ref{lem:at-sformulas} (ii) hold on every entry; the printed form before the revision (opposite sign) fails on every entry where it was tested.
\end{proposition}

\begin{proposition}[Two-leg solutions with all free directions]
\label{prop:at-c-sol}
$F_1^{(0)}|_2=\tfrac{27}{16}u\int(d^*\AV,d^*\AV)$ plus the $s$-closed arity-zero directions, and
\[
\begin{aligned}
\hbar F_1^{(1)}(\Delta)|_2&=u\Big[\tfrac{22}3\int(A,A^*)-\tfrac{47}3\int(\lambda^*,\lambda)-\tfrac83\int(c^*,c)\Big]+u\sum_d\theta^\Delta_dS(\Delta)_d\\
&\quad+\kappa_V\,uQ\!\int(\iota_x\AV,c^*)+\kappa_\lambda\,uQ\!\int(\rho(x)\lambda^*,\lambda^*),\\
\hbar F_1^{(1)}(P_b)|_2&=u\sum_d\theta^P_dS(P_b)_d+(\text{the }b\text{-parts of the }\kappa\text{ directions}),\\
\hbar F_1^{(1)}(\xi)|_2&=F^{\rm part}(\xi)+u\,\theta^\xi_VS^V_2(\xi)+u\,\theta^\xi_\perp S^\perp_2(\xi)+\{N_f,S_2(\xi)\}\text{-shifts},\\
\hbar F_1^{(1)}(\eta)|_2&=\sum_\mu x^\mu F^{\rm part}(\rho(e_\mu)\eta)-\big(3L-\tfrac32\big)u\,s_0-\tfrac52u\,u_0+u\,\theta^{\eta}_VS^V(\eta)+u\,\theta^{\eta}_\perp S^\perp(\eta),\\
\hbar F_1^{(1)}(M_\omega)|_2&=\tfrac13u\,S(M_\omega)_{\rm spin}+u\sum_d\theta^M_dS(M_\omega)_d,\qquad \hbar F_1^{(1)}(R_r)|_2=u\sum_d\theta^R_dS(R_r)_d,\\
\hbar F_1^{(1)}(K_b)|_2&=u\int(b,x)\Big\{-\tfrac{22}3(A,A^*)+16(\lambda,\lambda^*)+\tfrac83(c,c^*)\Big\}\\
&\quad+\tfrac23u\int\big\{(\iota_bA)(\iota_xA^*)-(\iota_xA)(\iota_bA^*)\big\}-\tfrac13u\int b_\mu x_\nu(\lambda,\rho_\mu\rho_\nu\lambda^*)+u\sum_d\theta^K_dS(K_b)_d+\sum_{j=1}^4\kappa^K_jG^K_j,
\end{aligned}
\]
with $Q$ of the first bracket equal to $\hbar\Ob^{(1)}(\Delta)|_2$ (gate G-F1Delta), $G^K_1,\dots,G^K_4$ the four second-order free-field symmetries of weight $3$ (closed, not exact, not of the form $S(z)$ on two legs; listed in action/complete/results.json, F1(K\_b)|2, free\_B), and the $x$-linear part of $\hbar F_1^{(1)}(K_b)|_2$ equal to $-(b,x)$ times the particular part along $\Delta$ plus the spin term of the rotation part of $d_{b'}v$ ($16-\tfrac13=\tfrac{47}3$). On two legs $H^0(Q)$ vanishes along $\Delta$, $P_b$, $M_\omega$, $R_r$, $\Sigma$, and is $1$-, $4$-dimensional along $\Sigma^*$ (position degree $\le1$), $K_b$ (position degree $\le2$); in both cases the residual has no component along it.
\end{proposition}

\begin{proposition}[Two-leg arity-two equations]
\label{prop:at-c-ar2}
On two legs, with the solutions of Proposition \ref{prop:at-c-sol}, the classes of the arity-two residuals vanish exactly under the following conditions (both routes):
\[
\begin{array}{l|l}
\text{pair} & \text{condition}\\
\hline
(\Delta,P_b) & \theta^\Delta_\perp=\theta^\Delta_\lambda=\theta^\Delta_V=-\tfrac13\quad(o_d=-\tfrac23u)\\
(\Delta,\xi_\pm) & \theta^\Delta_\perp=\theta^\Delta_\lambda=\theta^\Delta_V=-\tfrac13,\ \theta^\xi\ \text{free}\\
(\xi_+,\xi_-) & \theta^P_d=\theta^{\xi_+}_V+\theta^{\xi_-}_V-\tfrac{27}4-2L\ (\text{all }d),\quad \theta^{\xi_+}_V+\theta^{\xi_-}_V=\theta^{\xi_+}_\perp+\theta^{\xi_-}_\perp\\
(\xi_\pm,\xi'_\pm) & \theta^{\xi_\pm}_V=\theta^{\xi_\pm}_\perp\\
(\xi_\pm,P_b) & \theta^P_V=\theta^P_\perp=\theta^P_\lambda\\
(\xi_+,\eta_-) & \theta^\xi_V-\theta^\xi_\perp=\theta^\eta_V-\theta^\eta_\perp\\
(\xi_+,\eta_+) & \theta^R_\perp=\theta^R_\lambda=:\tau,\ \theta^\Delta_d=\theta^M_d=\tau-\tfrac13,\ \theta^\xi_d+\theta^\eta_d=\tau+\tfrac{77}{12}+2L\ (d=V,\perp)
\end{array}
\]
($\theta^\eta_d$ in the parametrisation of Proposition \ref{prop:at-c-sol}; in the parametrisation of action/complete/notes.txt 2.4, $\theta^\eta_\perp$ there $=\theta^\eta_\perp-3L+\tfrac32$ here.)
The obstructions: $\hbar\Ob^{(2)}(\Delta,\xi)|_2=u\big[8(T_1+T_2)-\tfrac{26}3(T_3+T_4)+\log\tfrac43\,Q(\psi_1+\psi_2)\big]$, $\psi_1=\int x^\mu(\Gamma(\xi,\lambda)^\mu,c^*)$, $\psi_2=\int x^\mu(A^*_\mu,(\xi,\lambda^*))$, solved by $\hbar F_1^{(2)}(\Delta,\xi)|_2=\log\tfrac43\,u(\psi_1+\psi_2)$ plus closed directions; $\hbar\Ob^{(2)}(\xi,P_b)|_2$ is $\tfrac14\log\tfrac43$ times a $Q$-exact functional; $\hbar\Ob^{(2)}(\Delta,P_b)|_2$ has no position part and no pole; every two-leg arity-two pole sum vanishes (along $(\xi,\eta)$ checked on the intrinsic part). With $\theta^\xi_V=\theta^\xi_\perp=:t_\xi$ the rows $(\xi_+,\eta_\pm)$ give $\theta^\eta_V=\theta^\eta_\perp=:t_\eta$ (the position part of $F_1^{(1)}(\eta)$ is then $x^\mu F_1^{(1)}(\rho(e_\mu)\eta)$ with one uniform coefficient) and uniform $t_P$, $t_\Delta$, $t_M$, $t_R$:
\[
t_\Delta=-\tfrac13,\qquad t_M=t_\Delta,\qquad t_R=t_\Delta+\tfrac13,\qquad t_{\xi_+}+t_{\xi_-}-t_P=\tfrac{27}4+2L,\qquad t_{\xi_+}+t_{\eta_+}-t_\Delta=\tfrac{27}4+2L .
\]
\end{proposition}

\begin{proof}
Computed: the cycles of Theorem \ref{thm:at-classes} and Proposition \ref{prop:at-classes2} in the two routes, the cohomology on complete bases fixed by representation counts (for $(\xi,\eta)$ decided on the value vectors directly, the generated $(\lambda,\lambda^*)$ position basis lacking one structure per chirality), and the residual formulas of Definition \ref{def:at-ob}. The arity-two generating function for two odd generators is $I+cS(\xi)+c'S(\xi')-cc'S^{(2)}(\xi,\xi')$, the sign forced by the classical master equation (gate GA-S2/GB-S2: $QS^{(2)}_2+S_2(P_{\Gamma|_V})-\{S_2(\xi),S_2(\xi')\}=0$). The last line: insert $\theta^\Delta=-\tfrac13$ into the $(\xi_+,\eta_+)$ row, giving $\tau=0$.
\end{proof}

\begin{corollary}[Scheme-independent content on the action side]
\label{cor:at-c-inv}
The inner direction ($t_z\mapsto t_z+c\lambda_z$, $\lambda=-1,+1,-2,+2,0$ on $\Sigma,\Sigma^*,P,K$ and the even rotations and $\Delta$) and the outer direction ($t_{\xi_\pm}\mapsto t_{\xi_\pm}\pm a$, $t_{\eta_\pm}\mapsto t_{\eta_\pm}\mp a$) leave invariant exactly the combinations
\[
t_\Delta=-\tfrac13,\qquad t_M-t_\Delta=0,\qquad t_R-t_\Delta=\tfrac13,\qquad t_{\xi_+}+t_{\xi_-}-t_P=\tfrac{27}4+2L,\qquad t_{\xi_+}+t_{\eta_+}-t_\Delta=\tfrac{27}4+2L,
\]
and the sector uniformity of $t_P$, $t_\xi$, $t_\eta$; these are the one-loop parts of $[D,P]=-2P$, $[D,Q]=-Q$, $[Q,P]=0$, $[Q,Q]=P$ and $\{Q,S\}=D+M+R$ in the form of Theorem \ref{thm:at-algebra}. In a translation-covariant scheme $t_P=t_\Delta=-\tfrac13$, hence $t_{\xi_+}+t_{\xi_-}=\tfrac{77}{12}+2L$ and $t_{\eta_+}=t_{\xi_-}$. Scheme artefacts: the individual $t_{\xi_\pm}$, $t_{\eta_\pm}$, $t_P$, $t_K$, $\theta^K_d$, $\kappa^K_j$, $\kappa_V$, $\kappa_\lambda$, the particular parts with $\log\tfrac32$, and every representative above.
\end{corollary}

\begin{remark}[Two-component forms]
\label{rem:at-c-tc}
With Remark \ref{rem:at-twocomp}: $\int(A,A^*)=\int d^4x\{(A_m,A^*_m)+\tfrac14(A_{ij},A^{*ij})\}$, $\int(\lambda^*,\lambda)=\int(\lambda,\lambda^*)=\int d^4x\{(\lambda_\alpha{}^i,\lambda^{*\alpha}{}_i)-(\bar\lambda_{\dot\alpha i},\bar\lambda^{*\dot\alpha i})\}$, $\int(\iota_bA)(\iota_xA^*)=\int d^4x\,(b_mA_m,x_nA^*_n)$, and $S^{(2)}(\xi,\xi')=\int\{\tfrac14(\Gamma(\xi,\xi'),\Gamma(\lambda^*,\lambda^*))-\tfrac12(\xi,\lambda^*)(\xi',\lambda^*)-(\iota_{\Gamma(\xi,\xi')}A,c^*)\}$ with $\Gamma(\xi,\bar\xi')|_V^m=i\,\bar\xi'_{\dot\alpha i}\bar\sigma^{m\dot\alpha\alpha}\xi_\alpha{}^i$, $\Gamma|_\perp=0$ for $\xi\in S_+\otimes W$, $\bar\xi'\in S_-\otimes W^*$, and $\Gamma|_V=0$, $(\Gamma|_\perp)_{kl}=2\varepsilon_{jikl}\,\xi_\alpha{}^j\varepsilon^{\alpha\beta}\xi'_\beta{}^i$ for $\xi,\xi'\in S_+\otimes W$. On the Dirac column the $\rho_\mu\rho_\nu$ of the spin terms is $\gamma^\mu_W\gamma^\nu_W$ (Proposition c8-twocomp), on the undotted block $\sigma^\mu\bar\sigma^\nu$ and on the dotted block $\bar\sigma^\mu\sigma^\nu$.
\end{remark}

\section{Engines and gates}
\label{sec:at-engines}

The specification is \texttt{engine\_spec.json} in this folder. Two independent routes compute every item; neither reads the other's outputs before the comparison.
\begin{itemize}
\item Route A (special-edge cycles in simplex form): enumerates the irreducible defect cycles of Theorems \ref{thm:at-classes} and Proposition \ref{prop:at-classes2} sector by sector, evaluates $a_m$ and $a_0$ exactly by Lemma \ref{lem:at-special} and Proposition \ref{prop:at-schwinger}, decomposes on exact bases of the target sectors by evaluation on integer plane waves in the explicit realisations of $V_{10}$, $\Sigma$ over $\mathbb Q(i)$ and $\mathfrak g=\mathfrak{su}(2)$, $\mathfrak{su}(3)$, $\mathfrak{gl}(3)$, and solves the equations by exact linear algebra. It extends the cycle enumerator of the Part-C action engine (step2/engine\_a, s2\_*) and the defect-chain code of Part Q (step2/composite\_qsk/qsk/engine\_a, qa\_core, qa\_s1) from chains through two letters to closed cycles.
\item Route B (renormalised amplitudes and $Q$): computes $W^{(1)}_{\rm ren}(x)$ and $W^{(2)}_{\rm ren}(x,y)$ by momentum-space numerator algebra with exact Feynman-parameter integrals, applies $Q$, reads the local part, and checks the special-edge cycles numerically at finite $\varepsilon$. It extends step2/engine\_b (ia\_engine, delta\_engine) and step2/composite\_qsk/qsk/engine\_b (qb\_core, qb\_numeps, qb\_numT).
\end{itemize}
Gates (regressions and identities of the construction): G-sign (the component formulas of Lemma \ref{lem:at-sformulas} and $s^2=0$ on random test cochains); G-reg0 (the arity-zero two-leg obstruction $-\tfrac{27}8u\int(d^*dc,d^*\AV)$ and the cubic obstruction \eqref{eq:at-ob3}); G-reg1 (\eqref{eq:at-reg2}); G-pole (\eqref{eq:at-pole} and \eqref{eq:at-local2}); G-log (no $\log\varepsilon$ in any cycle; $s\,\lc=0$); G-logex (Remark \ref{rem:at-logexact}, the closed combination $\lc^{(1)}(x)-\{uN,S(x)\}$); G-R0 ($\Ob^{(1)}(R_r)=0$ cycle sum); G-P2 ($\Ob^{(1)}(P_b)|_2=0$); G-pos (Proposition \ref{prop:at-pos}); G-WZ1, G-WZ2 (Lemma \ref{lem:at-sformulas} (iii) on two and three legs, on the full functional, intrinsic and position parts, with $\Lambda^{(0)}G^{(1)}(x)=-\{G^{(0)},S(x)\}$ and the arity-two form of (ii)); G-T (finite $T$ along $K_b$: $T\partial_T=O(\varepsilon)$); G-AB (route A against route B, entry by entry); G-tDelta (the determinations of $t_\Delta$ agree); G-zero (the classes on the commuting pairs vanish); G-H (Hypothesis \ref{hyp:at-H} at three legs; distinct from the two-leg cohomology gate G-S1xi2 of Proposition \ref{prop:at-xi2}). Cross-checks: X-scal, X-theta, X-twist, X-KP (Remark \ref{rem:at-cross}). Regression values:
\begin{equation}
\hbar\Ob^{(0)}\big|_3=u\Big[\big(\tfrac{11}4+\tfrac{19}4\log\tfrac32\big)(G_1+G_3)+\tfrac{27}8G_2-9\log\tfrac32\,G_5\Big]\quad(\text{Theorem c9-ob3}).
\label{eq:at-ob3}
\end{equation}

\section{Two-component forms}
\label{sec:at-twocomp}

\begin{remark}[Two-component forms of the objects of this note]
\label{rem:at-twocomp}
With Definition c8-twocomp ($\sigma^m=(\sigma^1,\sigma^2,\sigma^3,i)$, $\bar\sigma^m=(-\sigma^1,-\sigma^2,-\sigma^3,i)$, $\varepsilon^{12}=\varepsilon^{\dot1\dot2}=+1$), for $\xi\in S_+\otimes W$ with components $\xi_\alpha{}^i$ ($\bar\xi=0$), $\eta\in S_+\otimes W^*$ with components $\eta^\alpha{}_i$ ($\bar\eta=0$), and the pairing of $\mathfrak g$ written as $(-,-)$ around the components:
\[
\begin{aligned}
&\Gamma_V^m=\Gamma(\xi,\lambda)|_V^m=i\,\bar\lambda_{\dot\alpha i}\bar\sigma^{m\dot\alpha\alpha}\xi_\alpha{}^i,\qquad (\Gamma_\perp)_{kl}=2\varepsilon_{jikl}\,\xi_\alpha{}^j\varepsilon^{\alpha\beta}\lambda_\beta{}^i,\qquad (\xi,\lambda^*)=\xi_\alpha{}^i\lambda^{*\alpha}{}_i,\\
&T_0=-\int d^4x\,\big(\partial_mA_m,\ \xi_\alpha{}^i\lambda^{*\alpha}{}_i\big),\qquad T_1=-i\int d^4x\,\big(\bar\lambda_{\dot\alpha i}\bar\sigma^{m\dot\alpha\alpha}\xi_\alpha{}^i,A^*_m\big),\qquad T_2=-\tfrac12\int d^4x\,\varepsilon_{jikl}\big(\xi_\alpha{}^j\varepsilon^{\alpha\beta}\lambda_\beta{}^i,A^{*kl}\big),\\
&T_3=-\tfrac14\int d^4x\,\big((\partial_mA_n-\partial_nA_m)\,\sigma^m_{\alpha\dot\alpha}\bar\sigma^{n\dot\alpha\beta}\xi_\beta{}^i,\ \lambda^{*\alpha}{}_i\big),\qquad T_4=-\tfrac i2\int d^4x\,\big(\varepsilon^{\alpha\beta}\xi_\beta{}^j\sigma^m_{\alpha\dot\alpha}\partial_mA_{ij},\ \bar\lambda^{*\dot\alpha i}\big),\\
&g_1=\int d^4x\,\big(\Gamma_V^n,\ -\partial^2A_n+\partial_n\partial_mA_m\big),\quad g_2=-\int d^4x\,\big(\Gamma_V^n,\partial_n\partial_mA_m\big),\quad g_3=-\tfrac14\int d^4x\,\big((\Gamma_\perp)_{kl},\partial^2A^{kl}\big),\quad g_4=-\int d^4x\,\big(\partial^2c,\ \xi_\alpha{}^i\lambda^{*\alpha}{}_i\big),\\
&s_0=\int d^4x\,\big(A^{kj}\varepsilon_{\alpha\beta}\eta^\beta{}_j,\ \lambda^{*\alpha}{}_k\big),\qquad \int(\rho(\AV)\eta,\lambda^*)=-i\int d^4x\,\big(A_m\eta^\alpha{}_i\sigma^m_{\alpha\dot\alpha},\ \bar\lambda^{*\dot\alpha i}\big),
\end{aligned}
\]
with $A^{ij}=\tfrac12\varepsilon^{ijkl}A_{kl}$, $d^*=-\iota_\partial$, $d^*d\AV=-\partial^2A_n+\partial_n\partial_mA_m$ and $d^*d$ on the scalars equal to $-\partial^2$. The sector parts are
\[
\begin{aligned}
&S(\Delta)_\perp=\tfrac14\int d^4x\,\big((2x_n\partial_n+2)A_{ij},A^{*ij}\big),\qquad S(\Delta)_V=\int d^4x\,\Big\{\big((2x_n\partial_n+2)A_m,A^*_m\big)-\big(2x_n\partial_nc,c^*\big)\Big\},\\
&S(\Delta)_\lambda=-\int d^4x\,\Big\{\big((2x_n\partial_n+3)\lambda_\alpha{}^i,\lambda^{*\alpha}{}_i\big)-\big((2x_n\partial_n+3)\bar\lambda_{\dot\alpha i},\bar\lambda^{*\dot\alpha i}\big)\Big\},
\end{aligned}
\]
and $S(P_b)_d$ the same with $2x_n\partial_n+2w$ replaced by $b_n\partial_n$. The field-counting functional is $N=2\int(A,A^*)-4\int(\lambda^*,\lambda)-\int(c^*,c)$ with $\int(A,A^*)=\int d^4x\{(A_m,A^*_m)+\tfrac14(A_{ij},A^{*ij})\}$ and $\int(\lambda^*,\lambda)=\int d^4x\{(\lambda_\alpha{}^i,\lambda^{*\alpha}{}_i)-(\bar\lambda_{\dot\alpha i},\bar\lambda^{*\dot\alpha i})\}$. The regression \eqref{eq:at-reg2} reads
\[
\hbar\Ob^{(1)}(\Delta)\big|_2=u\int d^4x\Big\{-\tfrac{11}3\Big[\tfrac12(F^{\rm lin}_{mn},F^{\rm lin}_{mn})+\tfrac14(\partial_mA^{ij},\partial_mA_{ij})\Big]+\tfrac{94}3\,i\big(\bar\lambda_{\dot\alpha i},\bar\sigma^{m\dot\alpha\alpha}\partial_m\lambda_\alpha{}^i\big)-\tfrac{14}3\big(\partial_mc,A^*_m\big)\Big\},
\]
$F^{\rm lin}_{mn}=\partial_mA_n-\partial_nA_m$, and $G_1,\dots,G_5$ of \eqref{eq:at-ob3} are given in Remark c9-twocomp3. On the Dirac column $(\lambda_\alpha{}^i,\bar\lambda^{\dot\alpha}{}_i)$, $\rho(v)=\gamma_W^\mu v_\mu$ gives Weinberg's form of each line (Proposition c8-twocomp).
\end{remark}

\section{Status of the statements}
\label{sec:at-status}

\begin{itemize}
\item Derived from the construction (proofs above): Lemmas \ref{lem:at-defect}, \ref{lem:at-sformulas} (signs revised 2026-11-04 for the bracket of Definition c1-bracket), \ref{lem:at-place}, \ref{lem:at-special}; Propositions \ref{prop:at-poles}, \ref{prop:at-log}, \ref{prop:at-composite}, \ref{prop:at-pos}, \ref{prop:at-T}, \ref{prop:at-xi2}, \ref{prop:at-classes2}; Theorems \ref{thm:at-classes}, \ref{thm:at-second}, \ref{thm:at-bound}, \ref{thm:at-eq1}, \ref{thm:at-cubic}, \ref{thm:at-eq2} ((iii) by the forcing argument, not written out symbol by symbol); Theorems \ref{thm:at-H}, \ref{thm:at-fix}, \ref{thm:at-theta}, \ref{thm:at-algebra} under Hypothesis \ref{hyp:at-H}.
\item Computed by two engine routes and checked by the referee: Propositions \ref{prop:at-c-ob1}, \ref{prop:at-c-sol}, \ref{prop:at-c-ar2} and Corollary \ref{cor:at-c-inv} on two legs; the three-leg regressions \eqref{eq:at-ob3}; $\Ob^{(1)}(R_r)=0$ on three tested three-leg sectors.
\item Open (not computed): the pairs $(K_b,P_{b'})$, $(\eta,\eta')$, $(\xi_-,\eta_\pm)$; all three- and four-leg classes (Theorem \ref{thm:at-cubic}, task S1-cubic, Hypothesis \ref{hyp:at-H} at three legs, Proposition \ref{prop:at-couplings}); the four closed directions $G^K_j$ and the $\mathfrak{su}(4)$-$\mathbf{15}$ free classes; G-log as a finite-$\varepsilon$ extrapolation; G-T as a numerical finite-$T$ run; the five-leg sector along $P_b$.
\item Hypothesis: Hypothesis \ref{hyp:at-H} (classical symmetries at the listed weights), tested at three legs by G-H (not run).
\item Former shortcuts, now theorems from the construction or cross-checks: $\theta=-\tfrac13$ is computed from the construction twice ($(\Delta,P_b)$ and $(\Delta,\xi_\pm)$, Proposition \ref{prop:at-c-ar2}); the operator-level derivation (Remark \ref{rem:at-thetacheck}) is a cross-check only. The scaling identity ``$4$ times the logarithmic coefficient'' along $\Delta$ is Theorem \ref{thm:at-second} (v), used as the cross-check X-scal (it holds placement class by placement class). The twist (Remark \ref{rem:at-twist}) and $K=P^\dagger$ are cross-checks only and enter no number.
\end{itemize}

\end{document}
